\documentclass[11pt]{article}
\usepackage[utf8]{inputenc}
\usepackage[T1]{fontenc}
\usepackage[spanish,es-noquoting,es-tabla]{babel}
\usepackage{lmodern}
\usepackage[a4paper,margin=2.5cm]{geometry}

\setlength{\parindent}{0pt}
\setlength{\parskip}{8pt}

\begin{document}

\begin{center}
{\Large\bfseries Declaración de uso de IA generativa}\\[2pt]
Ciencia de Datos -- Unidad 2 -- Trabajo aplicativo -- Universidad de Cuenca
\end{center}

\textbf{Herramienta utilizada:} Claude (Anthropic), a través de Claude Code.

\textbf{Para qué la usamos:} el trabajo lo fuimos construyendo como grupo, discutiendo e iterando ideas sobre el problema, el actor, los indicadores y el diseño de la herramienta hasta llegar a la propuesta final. Usamos la IA como apoyo en ese proceso: la fuimos guiando con nuestras decisiones y correcciones para que nos ayudara a dar forma limpia y ordenada a los documentos finales (la matriz de casos, el Canvas, el diccionario de datos, el diseño, el plan de validación y el pitch) y a organizar el código del notebook.

\textbf{Cómo verificamos las salidas:} revisamos cada documento antes de subirlo al repositorio. Comprobamos que las cifras coincidieran con los resultados del notebook y con las fuentes que usamos (los boletines del proyecto LALA, el manual de FlipMyLearning y el estudio de Bravo, Peña e Illescas sobre reprobación y abandono en la UCuenca), y ejecutamos el notebook completo para confirmar que funciona.

El contenido final es responsabilidad del grupo. No compartimos datos personales de estudiantes de la Universidad de Cuenca; la prueba de concepto se hizo únicamente con el dataset público OULAD.

\end{document}
